\honest{Priming and Contamination}{Eliezer Yudkowsky}{2007}

Ask people to press one button if a string of letters is a word (``banner'') and another if it is not (``banack''). Show them ``water,'' and they will later recognize ``drink'' as a word a little faster. This is semantic priming. \nb{This kind of priming is well established.} It works at the level of recognizing that letters form a word, which one would expect to happen before thinking about meaning. It probably spreads to ``river,'' ``cup'' and ``splash'' too.

Priming ``is subconscious and unstoppable.'' \nb{In 1999 Stolz and Besner had argued against the view that reading activates meaning automatically: ``Semantic processing depends strongly on attentional control over how activation is distributed across different levels of representation.''} Try naming the ink colors of color words printed in the wrong colors, and you will see.

Anchors work partly this way. In Mussweiler and Strack's experiment, people asked whether Germany's mean temperature was above or below 5°C were then faster to recognize ``cold'' and ``snow''; people asked about 20°C were faster with ``hot'' and ``sun.'' Uninformative, false or irrelevant information affects estimates and decisions. This is called contamination.

Adjustment from an anchor still explains some cases, in particular when people produce the first estimate themselves, ``But modern research seems to show that most anchoring is actually due to contamination, not sliding adjustment.'' \nb{A month earlier, in ``Anchoring and Adjustment,'' Yudkowsky had called adjustment ``The current theory'' for the wheel experiment. The paper cited here, by Epley and Gilovich, describes an emerging consensus that no adjustment takes place, and argues that adjustment does occur when people produce the anchor themselves.}

Signs saying ``Limit 12 per customer'' or ``5 for \$10'' have ``been shown to work,'' which is why shops keep putting them up. You probably think you are not influenced, but someone must be.

Contamination, I conclude, is ``yet another of the thousand faces of confirmation bias.'' An idea in your head primes compatible thoughts ``and thereby ensures its continued existence.'' Quite apart from the pressure to win political arguments, confirmation bias is built into our hardware, in networks of association that bring up compatible thoughts and memories. ``All it takes is that one quick flash, and the bottom line is already decided.'' \nb{Every study in the post measures how fast words are recognized or what estimates people give in the same session. None measures whether a primed idea lasts. Several later findings that brief primes change judgments and behavior failed to replicate; the word priming at the start of the post has held up.}
